\chapter{Conclusion}
\label{ch:conclusion}

This thesis investigated the impact of preprocessing steps on the effectiveness of $z$-anonymity within multi-layered \ac{IoT} architectures. By analyzing the trade-offs between privacy, data utility, and system performance, this research aimed to identify the optimal architectural placement of data transformations in high-frequency smart metering systems.

The baseline evaluation confirmed that applying centralized $z$-anonymity to raw smart meter streams is not viable for operational use. Due to the high precision of the energy consumption data, individual readings are mathematically unique, which prevents the formation of sufficiently large anonymity sets. Consequently, enforcing strict privacy thresholds (e.g., $z=50$) without preprocessing leads to near-total data suppression, severely limiting data availability for the utility provider.

To address this limitation, a distributed architecture incorporating data preprocessing is recommended. Shifting the preprocessing workload to the Edge and Fog Layers proved highly effective in restoring data utility. Specifically, \textbf{Data Generalization} at the Edge emerged as the optimal strategy for balancing availability and precision. Moreover, the introduction of the $\text{NCP}_{\text{eff}}$ metric demonstrated that while coarse integer rounding ($p=0$) results in an excessive loss of granularity for typical households, rounding to two decimal places ($p=2$) provides an ideal compromise. It maintains a highly precise energy signal (yielding an $\text{NCP}_{\text{eff}}$ of only 1.2\%) while significantly improving the global Publication Ratio. 

Building upon this, combining Edge generalization with \textbf{Local Prefiltering} at the Fog Layer provides a comprehensive solution to balance privacy, utility, and performance. Implementing a low local threshold (e.g., $z_{\text{local}}=2$) protects data close to the source and successfully reduces the suppression workload at the Central Entity. This hybrid setup achieves approximately 44.0\% bandwidth savings. However, the local threshold must be carefully calibrated, as the restricted view of the population at the individual gateways can lead to the over-suppression of values that are otherwise globally common.

Finally, \textbf{Temporal Aggregation} proved to be an effective strategy strictly for optimizing system performance, delivering up to 87.5\% bandwidth savings. Nevertheless, the evaluation showed that averaging high-precision data generates new, unique arithmetic means, which prevents aggregation from improving the baseline Publication Ratio. Therefore, temporal aggregation is highly recommended for environments where network capacity is the primary constraint, but it must be paired with data generalization if a high volume of published tuples is required for operational analysis.

The widespread deployment of smart grids requires balancing operational data utility with consumer privacy. This thesis demonstrated that applying $z$-anonymity directly to raw, high-frequency smart meter data leads to excessive suppression due to the mathematical uniqueness of individual readings. However, integrating distributed preprocessing, specifically data generalization at the Edge Layer, effectively mitigates this issue by reducing data sparsity prior to the privacy check. Consequently, shifting the privacy-preserving workload to the data source is an essential architectural requirement for utilizing \ac{IoT} smart meter data securely and efficiently.

\newpage